\documentclass[12pt]{article}
\usepackage[utf8]{inputenc}
\usepackage[T1]{fontenc}
\usepackage{hyperref}
\title{Wa'ema Language Resources for K-12 Curriculum Development}
\author{Compiled by Claude Sonnet 5 (claude-sonnet-5) for Andr\'as Kornai}
\date{}
\begin{document}
\maketitle

\section{Introduction}
Wa'ema (Waema, Waiema; ISO 639-3 \texttt{wag}) is an Austronesian (Oceanic) language of Milne Bay Province, Papua New Guinea, spoken on the coast about 25 km west of Alotau. Glottolog places it in the Taupota dialect chain. Speaker counts vary, from about 1,000 in a 1997 dataset to 2,200 on Joshua Project. Very little educational material was found.

\section{Primary Educational Resources}
None were found.

\section{Linguistic and Reference Materials}
\begin{itemize}
\item Glottolog languoid page: \url{https://glottolog.org/resource/languoid/id/waem1237.bigmap.html}
\item Numeral system description (MPI-EVA): \url{https://lingweb.eva.mpg.de/channumerals/Wa%27ema.htm}
\item Taupota language (Wikipedia, related dialect chain): \url{https://en.wikipedia.org/wiki/Taupota_language}
\end{itemize}

\section{Religious and Community Resources}
Joshua Project reports no Christian resources in Wa'ema. It says the New Testament exists in two similar dialects.
\begin{itemize}
\item Joshua Project, Waema: \url{https://joshuaproject.net/people_groups/19562}
\item Joshua Project language page: \url{https://joshuaproject.net/languages/wag}
\end{itemize}

\section{Digital and Audio Resources}
None were found.

\section{Curriculum Development Context}
Wa'ema is reportedly spoken by all in the community, but no school programme was found. The related Taupota and Wedau varieties may offer comparative material. A curriculum would need new development with the community.
\end{document}